% Generated by tools/edn_latex.py from reports/edn.md.
% Do not edit: change reports/edn.md and run `make edn`.
\ednentry{
  id = {EDN-27},
  title = {Separate the requirements from their specification, keeping one set of FR/NFR IDs},
  date = {2026-09-29},
  milestone = {M1: Project Inception},
  topic = {Requirements Engineering},
  participants = {@lukas2510},
  decision = {\texttt{docs/\allowbreak{}docs/\allowbreak{}requirements.\allowbreak{}md} keeps the \texttt{FR-\allowbreak{}xx} and \texttt{NFR-\allowbreak{}xx} vocabulary but is rewritten at a design-free level: what must be true, for whom, with a priority and an acceptance criterion, and without endpoints, formats, status codes, libraries or deployment mechanics. Everything design-level moves to a new \texttt{docs/\allowbreak{}docs/\allowbreak{}specification.\allowbreak{}md}, which carries the same IDs, so \texttt{FR-\allowbreak{}04} there specifies \texttt{FR-\allowbreak{}04} here. Tests, the requirement-to-test matrix and all earlier EDN entries keep pointing at one unchanged ID space.},
  rationale = {\emph{Alternatives considered:}
\begin{itemize}[leftmargin=*, topsep=2pt]
\item \textbf{Option A (chosen): two levels, one ID space.}
\newline Pros: the requirements answer what the component must do and can be read by someone who does not know the architecture, which is what the feedback asked for; every design decision that was already analysed and settled (EDN-08, EDN-11 to EDN-18) survives in full, at the level where it belongs; a later design change, for example a different endpoint split or a different deployment path, no longer looks like a requirements change; the \texttt{FR-\allowbreak{}xx}/\texttt{NFR-\allowbreak{}xx} IDs that tests and twelve earlier EDN entries refer to keep working, and the requirement-to-test traceability is unaffected.
\newline Cons: two pages to keep consistent, and a reader has to know which of the two answers their question; every cross-reference in the brief, the problem specification, \texttt{data-\allowbreak{}versioning.\allowbreak{}md}, \texttt{CONTRIBUTING.\allowbreak{}md} and the EDN evidence links had to be repointed to the level it actually meant.
\item \textbf{Option B: rewrite the requirements as user stories (\texttt{US-\allowbreak{}xx}) with a specification behind them.}
\newline Pros: the most explicit way to show the user's intent, and the standard form in agile requirements engineering, which fits the Scrum setup of the course.
\newline Cons: introduces a second ID space on top of \texttt{FR-\allowbreak{}xx}/\texttt{NFR-\allowbreak{}xx}, so every test and every earlier EDN entry would reach its requirement through one more hop; the course material and the report template speak of functional and non-functional requirements, so the team would be renaming the artefact the report has to deliver. The team tried this form first and rejected it for these reasons; the abstraction it was meant to bring is achieved by the rewrite in option A without the renaming.
\item \textbf{Option C: keep one page and only soften its wording.}
\newline Pros: no new file, no link changes.
\newline Cons: the design decisions would still sit inside the requirements, only less visibly, so the criticism would be answered in form and not in substance; and there would be no place left for the endpoint-level detail that M4 needs.
\end{itemize}
\par
\emph{Why:} The lecturer's feedback on the M1 requirements was that they read as a specification: too technical, and fixing design where they should fix intent. The diagnosis holds. The page stated endpoints, HTTP status codes, header names, container mechanics and image tags, which meant that changing any of those looked like changing a requirement, and that a reader could not tell what the component owes its users from what we happened to choose. The detail itself is not the problem, since it is what the tests and the M4 deployment work are built on, so it was moved rather than dropped. Keeping a single ID space across both levels was the deciding factor against option B: the traceability the report has to show runs from a requirement ID to a test, and adding a second vocabulary on top of that would have cost a hop and gained nothing the rewrite does not already give.},
  ai = {Alternative generation, Alternative assessment, Recommendation, Solution generation},
  response = {Accepted with modifications},
  assessment = {Lukas relayed the feedback and asked for a more abstract version, with the explicit constraint that the existing analysis must not be lost. AI confirmed the criticism against the actual page and proposed the split into two levels, which the team accepted. Its first execution went further than asked and replaced the requirements with user stories; Lukas corrected this, because the course material, the report template and the existing IDs all speak of functional and non-functional requirements, and the stories had been meant as orientation only. AI then reworked the page to the agreed form. The abstraction of each requirement and the acceptance criteria are its proposal and were reviewed by the team.},
  aievidence = {Claude Code session on 2026-09-29 (prompts were in German, translated here): Lukas relayed the feedback as ``we got feedback on the requirements. The lecturer said they are too technical, they are more of a specification. He would rather have them like user stories, more abstract. Can you do that, but please also write a spec so the work was not for nothing'', and after AI delivered user stories corrected it with ``user stories were only meant as orientation for you. We still want to call them functional and non-functional requirements. But as said they should not be so technical, rather more abstract, and should not fix architecture decisions yet''.},
  evidence = {\href{https://github.com/mlops-2627q1-mds-upc/MLOPS_Recommenditos/blob/main/docs/docs/requirements.md}{Requirements}; \href{https://github.com/mlops-2627q1-mds-upc/MLOPS_Recommenditos/blob/main/docs/docs/specification.md}{Specification}; \href{https://github.com/mlops-2627q1-mds-upc/MLOPS_Recommenditos/pull/19}{PR \#19}.},
}
